\NormalFont\ShortTitle{1 Ioan}

{\MT Prima scrisoare a lui Ioan


\par }\ChapOne{1}{\PP \VerseOne{1}Ceea ce a fost de la început, ceea ce am auzit, ceea ce am văzut cu ochii noștri, ceea ce am văzut și am atins cu mâinile noastre, despre Cuvântul vieții

\VS{2}(și viața a fost descoperită, și noi am văzut, și mărturisim și vă vestim viața, viața veșnică, care era la Tatăl și care ne-a fost descoperită);

\VS{3}ceea ce am văzut și am auzit, vă vestim vouă, ca și voi să aveți părtășie cu noi. Da, iar părtășia noastră este cu Tatăl și cu Fiul Său, Isus Hristos.

\VS{4}Și vă scriem aceste lucruri, pentru ca bucuria noastră să fie împlinită.


\par }{\PP \VS{5}Iată mesajul pe care l-am auzit de la El și pe care vi-l vestim: Dumnezeu este lumină și în El nu este întuneric deloc.

\VS{6}Dacă spunem că avem părtășie cu El și umblăm în întuneric, mințim și nu spunem adevărul.

\VS{7}Dar dacă umblăm în lumină, așa cum și el este în lumină, avem părtășie unii cu alții, iar sângele lui Isus Hristos, Fiul său, ne curăță de orice păcat.

\VS{8}Dacă spunem că nu avem niciun păcat, ne înșelăm pe noi înșine și adevărul nu este în noi.

\VS{9}Dacă ne mărturisim păcatele, el este credincios și drept să ne ierte păcatele și să ne curățească de orice nelegiuire.

\VS{10}Dacă spunem că nu am păcătuit, îl facem mincinos și cuvântul Lui nu este în noi.



\par }\Chap{2}{\PP \VerseOne{1}Copilașii mei, vă scriu aceste lucruri ca să nu păcătuiți. Dacă cineva păcătuiește, avem un sfătuitor la Tatăl, pe Isus Hristos, cel neprihănit.

\VS{2}El este jertfa de ispășire pentru păcatele noastre, și nu numai pentru ale noastre, ci și pentru cele ale întregii lumi.

\VS{3}Iată cum știm că îl cunoaștem: dacă păzim poruncile Lui.

\VS{4}Cine spune: “Îl cunosc” și nu-i păzește poruncile este un mincinos și adevărul nu este în el.

\VS{5}Dar dragostea lui Dumnezeu a fost cu siguranță desăvârșită în oricine îi păzește cuvântul. Iată cum cunoaștem că suntem în el:

\VS{6}cine spune că rămâne în el trebuie să umble și el însuși așa cum a umblat el.


\par }{\PP \VS{7}Fraților, nu vă scriu o poruncă nouă, ci o poruncă veche, pe care ați avut-o de la început. Vechea poruncă este cuvântul pe care l-ați auzit de la început.

\VS{8}Vă scriu din nou o poruncă nouă, care este adevărată în el și în voi, pentru că întunericul trece și lumina cea adevărată strălucește deja.

\VS{9}Cel care spune că este în lumină și urăște pe fratele său este în întuneric până acum.

\VS{10}Cine iubește pe fratele său rămâne în lumină și nu există în el niciun prilej de poticnire.

\VS{11}Dar cel care își urăște fratele este în întuneric, umblă în întuneric și nu știe încotro se îndreaptă, pentru că întunericul i-a orbit ochii.


\par }{\PP \VS{12}Copilașilor, vă scriu vouă, pentru că păcatele v-au fost iertate pentru Numele Lui.


\par }{\PP \VS{13}Vă scriu vouă, părinților, pentru că voi cunoașteți pe Cel ce este de la început.

\par }{\PP Vă scriu vouă, tinerilor, pentru că ați biruit pe cel rău.

\par }{\PP Vă scriu vouă, copilașilor, pentru că voi îl cunoașteți pe Tatăl.


\par }{\PP \VS{14}V-am scris vouă, părinților, pentru că voi cunoașteți pe Cel ce este de la început.

\par }{\PP V-am scris vouă, tineri, pentru că sunteți tari, pentru că cuvântul lui Dumnezeu rămâne în voi și l-ați biruit pe cel rău.


\par }{\PP \VS{15}Să nu iubiți lumea și lucrurile din lume. Dacă cineva iubește lumea, dragostea Tatălui nu este în el.

\VS{16}Căci tot ce este în lume — pofta cărnii, pofta ochilor și mândria vieții — nu este al Tatălui, ci al lumii.

\VS{17}Lumea trece cu poftele ei, dar cel care face voia lui Dumnezeu rămâne pentru totdeauna.


\par }{\PP \VS{18}Copilașilor, acestea sunt vremurile de pe urmă și, după cum ați auzit că va veni Antihristul, chiar acum s-au ridicat mulți antihriști. Prin aceasta știm că este ceasul de pe urmă.

\VS{19}Ei au ieșit de la noi, dar nu ne aparțineau; căci, dacă ne-ar fi aparținut, ar fi continuat cu noi. Dar au plecat, ca să se arate că niciunul dintre ei nu ne aparține.

\VS{20}Voi aveți o ungere de la Cel Sfânt și toți aveți cunoștință.

\VS{21}Nu v-am scris pentru că nu cunoașteți adevărul, ci pentru că îl cunoașteți și pentru că nicio minciună nu este din adevăr.

\VS{22}Cine este mincinosul, dacă nu cel care neagă că Isus este Hristosul? Acesta este Antihristul, cel care îi neagă pe Tatăl și pe Fiul.

\VS{23}Cine îl neagă pe Fiul nu îl are pe Tatăl. Cel care Îl mărturisește pe Fiul Îl are și pe Tatăl.


\par }{\PP \VS{24}De aceea, să rămână în voi ceea ce ați auzit de la început. Dacă ceea ce ați auzit de la început rămâne în voi, și voi veți rămâne în Fiul și în Tatăl.

\VS{25}Aceasta este făgăduința pe care ne-a promis-o: viața veșnică.


\par }{\PP \VS{26}V-am scris aceste lucruri despre cei ce vor să vă ducă în rătăcire.

\VS{27}Cât despre voi, ungerea pe care ați primit-o de la el rămâne în voi și nu aveți nevoie ca cineva să vă învețe. Ci, așa cum ungerea lui vă învață despre toate lucrurile, și este adevărată și nu este minciună, și chiar așa cum v-a învățat el, veți rămâne în el.


\par }{\PP \VS{28}Și acum, copilașilor, rămâneți în El, pentru ca, atunci când Se va arăta, să avem îndrăzneală și să nu fim rușinați înaintea Lui, la venirea Lui.

\VS{29}Dacă știți că El este drept, știți că oricine practică dreptatea s-a născut din El.



\par }\Chap{3}{\PP \VerseOne{1}Vedeți cât de mare este dragostea pe care ne-a dat-o Tatăl, ca să ne numim copii ai lui Dumnezeu! De aceea, lumea nu ne cunoaște, pentru că nu l-a cunoscut pe el.

\VS{2}Iubiților, acum suntem copii ai lui Dumnezeu. Încă nu ni s-a descoperit ce vom fi; dar știm că, atunci când se va descoperi, vom fi asemenea lui, pentru că îl vom vedea așa cum este el.

\VS{3}Oricine are această nădejde pusă în el se purifică pe sine, așa cum și el este curat.


\par }{\PP \VS{4}Oricine păcătuiește comite și fărădelege. Păcatul este fărădelege.

\VS{5}Știți că El S-a arătat ca să ia păcatele noastre, și niciun păcat nu este în El.

\VS{6}Oricine rămâne în El nu păcătuiește. Cine păcătuiește nu l-a văzut și nu-l cunoaște.


\par }{\PP \VS{7}Copilașilor, nimeni să nu vă abată din drum. Cel ce face dreptate este drept, așa cum este el drept.

\VS{8}Cel ce păcătuiește este de la diavol, căci diavolul păcătuiește de la început. În acest scop a fost revelat Fiul lui Dumnezeu: ca să distrugă lucrările diavolului.

\VS{9}Oricine este născut din Dumnezeu nu comite păcat, pentru că sămânța lui rămâne în el și nu poate păcătui, pentru că este născut din Dumnezeu.

\VS{10}În aceasta s-au arătat copiii lui Dumnezeu și copiii diavolului. Oricine nu face dreptate nu este de la Dumnezeu, și nici cel care nu-și iubește fratele.

\VS{11}Căci acesta este mesajul pe care l-ați auzit de la început, că trebuie să ne iubim unii pe alții —

\VS{12}spre deosebire de Cain, care a fost de la cel rău și și-a ucis fratele. De ce l-a ucis? Pentru că faptele lui erau rele, iar ale fratelui său drepte.


\par }{\PP \VS{13}Nu vă mirați, frații mei, dacă lumea vă urăște.

\VS{14}Noi știm că am trecut din moarte la viață, pentru că iubim pe frații noștri. Cel care nu-și iubește fratele rămâne în moarte.

\VS{15}Oricine își urăște fratele este un ucigaș, și știți că niciun ucigaș nu are viață veșnică rămasă în el.


\par }{\PP \VS{16}Întru aceasta cunoaștem dragostea, pentru că El Și-a dat viața pentru noi. Și noi trebuie să ne dăm viața pentru frați.

\VS{17}Dar oricine are bunurile lumii și vede pe fratele său în nevoie, apoi își închide inima de compasiune împotriva lui, cum rămâne în el dragostea lui Dumnezeu?


\par }{\PP \VS{18}Copilașii mei, să nu iubim numai cu vorba și numai cu limba, ci cu fapta și cu adevărul.

\VS{19}Și prin aceasta cunoaștem că suntem din adevăr și ne convingem inimile înaintea lui,

\VS{20}pentru că, dacă inima noastră ne osândește, Dumnezeu este mai mare decât inima noastră și cunoaște toate lucrurile.

\VS{21}Preaiubiților, dacă inima noastră nu ne condamnă, avem îndrăzneală față de Dumnezeu;

\VS{22}așa că orice cerem, primim de la el, pentru că păzim poruncile lui și facem ceea ce este plăcut înaintea lui.

\VS{23}Aceasta este porunca lui: să credem în numele Fiului său, Isus Hristos, și să ne iubim unii pe alții, așa cum a poruncit el.

\VS{24}Cine păzește poruncile Lui rămâne în El și El în El. Prin aceasta cunoaștem că el rămâne în noi, prin Duhul pe care ni l-a dat.



\par }\Chap{4}{\PP \VerseOne{1}Preaiubiților, nu credeți orice duh, ci încercați duhurile, ca să vedeți dacă sunt de la Dumnezeu, pentru că mulți prooroci mincinoși au ieșit în lume.

\VS{2}Întru aceasta cunoașteți Duhul lui Dumnezeu: orice duh care mărturisește că Isus Cristos a venit în trup este de la Dumnezeu,

\VS{3}iar orice duh care nu mărturisește că Isus Cristos a venit în trup nu este de la Dumnezeu; și acesta este duhul lui Antihrist, despre care ați auzit că vine. Acum este deja în lume.

\VS{4}Voi, copilașilor, sunteți de la Dumnezeu și i-ați biruit, pentru că mai mare este Cel care este în voi decât cel care este în lume.

\VS{5}Ei sunt din lume. De aceea ei vorbesc din lume și lumea îi ascultă.

\VS{6}Noi suntem de la Dumnezeu. Cine îl cunoaște pe Dumnezeu ne ascultă. Cel care nu este din Dumnezeu nu ne ascultă. Prin aceasta cunoaștem duhul adevărului și duhul rătăcirii.


\par }{\PP \VS{7}Iubiților, să ne iubim unii pe alții, căci dragostea este de la Dumnezeu; și oricine iubește s-a născut din Dumnezeu și cunoaște pe Dumnezeu.

\VS{8}Cine nu iubește nu-L cunoaște pe Dumnezeu, căci Dumnezeu este iubire.

\VS{9}Prin aceasta s-a arătat în noi dragostea lui Dumnezeu: Dumnezeu l-a trimis pe singurul său Fiu născut în lume, pentru ca noi să trăim prin el.

\VS{10}În aceasta constă dragostea: nu că noi l-am iubit pe Dumnezeu, ci că el ne-a iubit pe noi și l-a trimis pe Fiul său ca jertfă de ispășire pentru păcatele noastre.

\VS{11}Iubiților, dacă Dumnezeu ne-a iubit în felul acesta, și noi trebuie să ne iubim unii pe alții.

\VS{12}Nimeni nu L-a văzut pe Dumnezeu în niciun moment. Dacă ne iubim unii pe alții, Dumnezeu rămâne în noi, iar dragostea Lui a fost desăvârșită în noi.


\par }{\PP \VS{13}Întru aceasta cunoaștem că noi rămânem în El și El în noi, pentru că ne-a dat din Duhul Său.

\VS{14}Noi am văzut și mărturisim că Tatăl a trimis pe Fiul ca Mântuitor al lumii.

\VS{15}Oricine mărturisește că Isus este Fiul lui Dumnezeu, Dumnezeu rămâne în el și el în Dumnezeu.

\VS{16}Noi cunoaștem și am crezut în dragostea pe care Dumnezeu o are pentru noi. Dumnezeu este iubire, iar cel care rămâne în iubire rămâne în Dumnezeu, iar Dumnezeu rămâne în el.

\VS{17}În aceasta, dragostea a fost desăvârșită între noi, ca să avem îndrăzneală în ziua judecății, pentru că, așa cum este el, așa suntem și noi în această lume.

\VS{18}În dragoste nu este frică, ci dragostea desăvârșită izgonește frica, pentru că frica are pedeapsă. Cel care se teme nu este desăvârșit în dragoste.

\VS{19}Noi îl iubim, pentru că el ne-a iubit primul.

\VS{20}Dacă un om spune: “Îl iubesc pe Dumnezeu” și îl urăște pe fratele său, este un mincinos; căci cine nu-l iubește pe fratele său, pe care l-a văzut, cum poate să-L iubească pe Dumnezeu, pe care nu L-a văzut?

\VS{21}Această poruncă o avem de la el: cine iubește pe Dumnezeu să iubească și pe fratele său.



\par }\Chap{5}{\PP \VerseOne{1}Oricine crede că Isus este Hristosul s-a născut din Dumnezeu. Cine iubește pe Tatăl iubește și pe copilul născut din el.

\VS{2}Prin aceasta cunoaștem că iubim pe copiii lui Dumnezeu, când îl iubim pe Dumnezeu și păzim poruncile lui.

\VS{3}Căci aceasta este iubirea față de Dumnezeu: să păzim poruncile lui. Poruncile Lui nu sunt grele.

\VS{4}Căci tot ce este născut din Dumnezeu biruiește lumea. Aceasta este victoria care a biruit lumea: credința voastră.

\VS{5}Cine este cel care biruiește lumea, dacă nu cel care crede că Isus este Fiul lui Dumnezeu?


\par }{\PP \VS{6}Acesta este Cel ce a venit prin apă și sânge, Isus Hristos, nu numai cu apă, ci cu apă și cu sânge. Duhul este cel care mărturisește, pentru că Duhul este adevărul.

\VS{7}Căci trei sunt cei care mărturisesc:

\VS{8}Duhul, apa și sângele; și cei trei sunt de acord ca unul singur.

\VS{9}Dacă noi primim mărturia oamenilor, mărturia lui Dumnezeu este mai mare, căci aceasta este mărturia lui Dumnezeu, pe care a mărturisit-o despre Fiul său.

\VS{10}Cel care crede în Fiul lui Dumnezeu are mărturia în el însuși. Cel care nu crede în Dumnezeu s-a făcut mincinos, pentru că nu a crezut în mărturia pe care Dumnezeu a dat-o despre Fiul său.

\VS{11}Mărturia este aceasta: Dumnezeu ne-a dat viața veșnică, iar această viață este în Fiul său.

\VS{12}Cine Îl are pe Fiul are viața. Cel care nu Îl are pe Fiul lui Dumnezeu nu are viața.


\par }{\PP \VS{13}Am scris aceste lucruri pentru voi, care credeți în Numele Fiului lui Dumnezeu, ca să știți că aveți viața veșnică și să credeți în continuare în Numele Fiului lui Dumnezeu.


\par }{\PP \VS{14}Îndrăzneala pe care o avem față de El este că, dacă cerem ceva după voia Lui, El ne ascultă.

\VS{15}Și dacă știm că ne ascultă, orice am cere, știm că avem cererile pe care i le-am adresat.


\par }{\PP \VS{16}Dacă cineva vede pe fratele său săvârșind un păcat care nu duce la moarte, să ceară și Dumnezeu îi va da viață pentru cei ce nu fac păcate care duc la moarte. Există păcat care duce la moarte. Nu spun că ar trebui să facă o cerere cu privire la aceasta.

\VS{17}Orice nelegiuire este păcat și există păcat care nu duce la moarte.


\par }{\PP \VS{18}Știm că oricine este născut din Dumnezeu nu păcătuiește; dar cel născut din Dumnezeu se păzește pe sine însuși, și cel rău nu se atinge de el.

\VS{19}Știm că noi suntem de la Dumnezeu, iar întreaga lume se află în puterea celui rău.

\VS{20}Știm că Fiul lui Dumnezeu a venit și ne-a dat pricepere, ca să cunoaștem pe cel adevărat; și noi suntem în cel adevărat, în Fiul său Isus Hristos. Acesta este adevăratul Dumnezeu și viața veșnică.


\par }{\PP \VS{21}Copilașilor, feriți-vă de idoli.


\par }